\begin{abstract}
Carbon dioxide removal is increasingly seen as a necessary addition to emission reductions in pathways consistent with the goals of the Paris Agreement. However, existing carbon markets remain dominated by low-quality carbon offsets, while the role of durable and scalable negative emission technologies (NETs) is still insufficiently understood. This thesis develops a dynamic portfolio model for four selected NETs: afforestation and reforestation (A/R), biomass carbon removal and storage (BiCRS), enhanced rock weathering (ERW), and direct air capture (DAC). Using bottom-up techno-economic models with country-level resolution, it estimates availability and marginal abatement costs for each approach and integrates them into static and dynamic optimization frameworks. The results show that A/R is the cheapest option, but its potential is constrained by land availability and climate conditions. BiCRS serves as an important intermediate technology, while DAC becomes crucial in the second half of the century as cheaper options reach their limits. Enhanced weathering exhibits substantial theoretical potential at competitive costs, but its practical contribution is limited by logistics and acceptance. Overall, the thesis demonstrates that no single technology can satisfy future CDR demand alone. Instead, a diversified portfolio of NETs is required, combining near-term natural sinks with engineered ("novel") approaches. The findings have important implications for investors and policymakers seeking to align carbon removal development and deployment with long-term climate targets.
\end{abstract}